\section{Completing the exponent-six census}
\label{sec:zero-test}
\textbf{Theorem.} The 432 fields in the exponent-six row of Elsenhans, Kl\"uners and Nicolae's Table~1 are all the imaginary quadratic number fields of exact class-group exponent six. Among them the largest absolute discriminant is $5\,761\,140$. A quadratic class group has exponent six when its sixth-power map is identically one and six is the least positive exponent with this property.

The finite enumeration of Elsenhans, Kl\"uners and Nicolae (EKN) is complete for $|D|\le X_0=3.1\cdot10^{20}$. Their Table~2, assuming the quadratic case of Dirichlet GRH, bounds the discriminant of every exponent-six field by $X_1=2.5\cdot10^{25}$. Indeed the analytic input is a bound for a rational prime splitting in $\mathbb Q(\sqrt D)$, whose zeta function is $\zeta(s)L(s,\chi_D)$. These are exactly the two zero statements available from the repository. We show that an exponent-six field with $X_0<|D|\le X_1$ cannot exist.
\href{https://arxiv.org/pdf/1803.02056}{Elsenhans, Kl\"uners and Nicolae (EKN), Theorems 1--3 and Tables 1--2};
\href{https://doi.org/10.1090/S0025-5718-96-00763-6}{Bach and Sorenson, Table~3}.

Write $x=|D|$, $e$ for the class-group exponent and
$\chi(n)=(D/n)$ for the primitive odd quadratic character of conductor $x$.
If $p$ splits and $P$ is either prime above $p$, then $P^e$ is principal.
Its generator $(a+b\sqrt D)/2$ has $b\ne0$ because its valuations at $P$
and $\bar P$ differ. Norms give $x\le 4p^e$. Consequently, if
$x>4Y^e$, every prime $p\le Y$ is inert or ramified. The explicit
formula can distinguish this pattern from the one required by the zeros
of $L(s,\chi)$.

\medskip\noindent\textbf{An explicit-formula inequality.}
Let $g$ be a compactly supported, real step function with nonzero norm.
Set
\[
 F(u)=\frac{\int_{\mathbb R}g(v)g(v+u)\,dv}{\int_{\mathbb R}g(v)^2\,dv},
 \qquad h(t)=\int_{\mathbb R} F(u)e^{itu}\,du.
\]
Then $F$ is even, continuous, piecewise linear, with $F(0)=1$ and
$|F(u)|\le1$. Its Fourier transform $h$ is nonnegative on the real line,
being a squared modulus there. It extends to an entire function with
$O((1+|t|)^{-2})$ decay on every fixed horizontal strip, by the explicit
Fourier integral of a step function and its reflected copy. Write $C$
for the largest absolute slope of $F$. For $\psi=\Gamma'/\Gamma$ the
Guinand--Weil formula is
\[
 \sum_{\rho}h\big((\rho-1/2)/i\big)=
 \log(x/\pi)+\frac{1}{2\pi}\int_{\mathbb R}h(t)
        \Re\psi(3/4+it/2)\,dt
 -2\sum_{n\ge2}\frac{\chi(n)\Lambda(n)}{\sqrt n}F(\log n).
\]
The left side is nonnegative by the repository's Dirichlet GRH. This is
where the zeros enter the bound. The odd parity determines the argument
$3/4$ of the digamma function.
\href{https://arxiv.org/pdf/2404.08380}{Carneiro, Milinovich,
Quesada-Herrera and Ramos, Lemma~9, with Fourier conventions converted}.

By the integral formula for the digamma function and Fourier inversion,
the integral including its factor $1/(2\pi)$ equals
\[
 \psi(3/4)+
 2\int_0^\infty \frac{e^{-3u/2}}{1-e^{-2u}}(1-F(u))\,du
 \ \le\ \psi(3/4)+\frac{14}9 C .
\]
Indeed the difference in the digamma integrand involves $1-\cos(tu)$;
its nonnegativity against $h(t)$ permits interchange of integrals.
Use $1-F(u)\le C u$, expand the denominator geometrically, and bound
$\sum_{j\ge0}(2j+3/2)^{-2}\le4/9+1/3$ by the integral test.
Also $\log\pi-\psi(3/4)>2$. To see this, $\log\pi>1.1$ and
$\gamma>H_6-\log 7>1/2$, while the first two terms of the digamma
difference give $\psi(1)-\psi(3/4)>1/3+1/14$ and $\psi(1)=-\gamma$.
It follows that
\[
 \log x> 2\sum_{n\ge2}\frac{\chi(n)\Lambda(n)}{\sqrt n}
                    F(\log n)+2-\frac{14}9 C.
\]

\medskip\noindent\textbf{A finite split-prime test.}
For each support bound $Y$ take $g$ constant on twelve successive intervals of width
$\delta=\log(Y)/12$, beginning at zero, with respective values
\[
          21,\ 14,\ 10,\ 6,\ 4,\ 1,\ -1,\ -4,\ -6,\ -10,\ -14,\ -21
\]
and zero elsewhere. Thus $F$ vanishes for $|u|\ge\log Y$. The opposite
signs at the ends allow inert primes to give a positive term in the
explicit formula even though $\chi(p)=-1$. Define
\[
 R_p=2\sum_{p^k\le Y}\frac{(-1)^k}{\sqrt{p^k}}F(k\log p),\qquad
 S=\sum_{p\le Y}R_p\log p,\qquad B=\max(0,\max_{p\le Y}R_p).
\]
If no $p\le Y$ splits, the prime sum including its factor two is
$S-\sum_{p\mid D,\,p\le Y}R_p\log p\ge S-B\log x$.
Hence that splitting pattern forces
\[
                 \log x>\frac{S+2-14C/9}{1+B}.
\]

\subsection{Forcing two equal coefficients}
Assume $e=6$ and $X_0<x\le X_1$. Because $4\cdot2000^6<X_0$, all split primes exceed 2000. Take $Y=5500$. Let $N_Y$ be the number of split primes up to $Y$. For these primes only $k=1$ occurs in the explicit formula. Changing a prime from inert to split subtracts the quantity
\[
   L_p=-4(\log p) F(\log p)/\sqrt p
\]
from the inert-model sum. The contribution lost at all ramified primes is at most $B\log x$, since $\sum_{p\mid D}\log p\le\log x$. Thus, for any common upper bound $L>0$ for the split-prime losses,
\[
 N_Y L > S+2-14C/9-(1+B)\log X_1.
\]
The exact rational enclosures in \texttt{zero\_test.py} give the two strict inequalities
\[
 \max_{2000<p\le5500} L_p<0.216,\qquad
 S+2-14C/9-(1+B)\log X_1>5.84.
\]
For orientation $S$ is about $65.682$, $B$ about $0.04476$, and $C$ about $0.47355$. Since $27\cdot0.216=5.832<5.84$, there are at least 28 split primes. Every such prime gives a norm equation
\[
       a_p^2+x b_p^2=4p^6,\qquad a_p\ge0,\quad b_p\ge1 .
\]
Here one may take absolute values of the two coefficients. The exact inequality $4\cdot5500^6<19^2 X_0$ gives $b_p\le18$. By the pigeonhole principle two distinct primes $p<q$ share the same $b$.

\subsection{Eliminating the discriminant}
Write $A=a_p$, $B'=a_q$. These nonnegative integers have the same parity, since their squares differ by a multiple of four. Subtract the two norm equations and put
\[
   u=(B'-A)/2,\quad v=(B'+A)/2,\quad
   N=q^6-p^6=uv.
\]
We have $0<u\le v$, $A=v-u$, and the positive integer
\[
                M=4p^6-(v-u)^2=x b^2
\]
is greater than $X_0$. Since $N=(q^3-p^3)(q^3+p^3)$ and $N/u-u$ decreases for $u>0$, positivity of $M$ is equivalent, for $u\le\sqrt N$, to $u>q^3-p^3$. Consequently it suffices to examine the positive divisors in
\[
       q^3-p^3<u\le\lfloor\sqrt{q^6-p^6}\rfloor .
\]
For an odd prime $r\le2000$, a nonzero square $-M\bmod r$ forces $r\nmid bD$ and $(D/r)=1$ because $-M=D b^2$. Such an $r$ would split in the candidate field. This contradicts $x>4\cdot2000^6$. It suffices to find one such witness for every remaining $M$; no extraction of its squarefree part is needed.

\paragraph{Finite lemma.} For each pair of primes $p<q\le5500$ satisfying $4p^6>X_0$, every divisor in the displayed interval yielding $M>X_0$ has a witness $r$ as above. The following counts include multiplicities over prime pairs and divisors.
\begin{center}
\begin{tabular}{lr}\toprule
Quantity & Count\\\midrule
Primes $p\le5500$ with $4p^6>X_0$ &414\\
Unordered pairs &85,491\\
Divisors in the strict-lower-endpoint interval &10,144,433\\
Pairs and divisors giving $M>X_0$ &9,393,302\\
Candidates with no witness &0\\\bottomrule
\end{tabular}
\end{center}
\textit{Verification.} Factor
\[
 N=(q-p)(q+p)(q^2+pq+p^2)(q^2-pq+p^2).
\]
Trial division is sufficient because each factor is at most $3\cdot5500^2$. Generate the divisors from their prime exponents, test the two inequalities and reduce $-M$ modulo the small primes. The Python implementation uses arbitrary-precision integers and sets of nonzero squares. The separately written C++ implementation uses 128-bit products, checks its factorization, and tests the equivalent residue $(v-u)^2-4p^6$ via precomputed masks. It includes the lower endpoint, which adds one divisor per pair with $M=0$ and affects none of the positive candidates. In this run primes from 19 to 151 already supplied a witness in every case. Both programs returned 9,393,302 tests above the threshold and zero survivors. The finite lemma contradicts the forced shared coefficient, proving the theorem.

\paragraph{The 432 fields and the order classification.}
The file \texttt{data/exponent-six.csv} lists every field discriminant and class number. For two through seven distinct ramified primes the counts in order are $60,141,140,76,12,3$. Aiken, Ross, Shvydkoy and Xue's Theorem~1.3(5) states that under ERH their table of exponent-six negative order discriminants is complete if EKN's remaining interval contains no fundamental example. Their construction reduces all analytic hypotheses to the underlying quadratic classifications. The integral endpoint $X_0$ is already covered by EKN (and is not a fundamental discriminant). Our theorem discharges the remaining condition. Thus exactly \textbf{1,236} imaginary quadratic orders have proper ideal class group of exponent six, of which $1,236-432=804$ are nonmaximal. The largest absolute order discriminant is $43\,522\,752$. We checked the 1,236 distinct entries in the authors' data and verified that its subset of 432 fundamental discriminants agrees entry for entry with the field table.
\href{https://github.com/erick-ross/research-data/blob/119301b8bcae42a891062d4877b77b6dbc8e820f/class-groups-small-exp/Table-6.txt}{Pinned order data}.
\href{https://arxiv.org/pdf/2605.26020}{Aiken, Ross, Shvydkoy and Xue, Theorem~1.3(5), Section~6 and Table~A.6}.
Using primitive positive reduced binary quadratic forms we verified the listed group orders and the sixth power of every one of the $103\,830$ classes in the 1,236 orders. More explicitly, after changing variables so that the leading coefficient $a$ is prime to the order discriminant, the form represents an invertible ideal $(a,\omega-u)$ in a quadratic order $\mathbb Z[\omega]$. Hensel-lifting the simple root $u$ of the polynomial of $\omega$ modulo $a$ gives $u_6$ modulo $a^6$. The ideal $(a^6,\omega-u_6)$ equals the sixth power by containment and equality of norms. Reduction of its norm form checks principality. Counting the reduced forms gives the order of the class group; a finite abelian group killed by six has elementary abelian Sylow subgroups, so its order determines its structure. This independently checks membership throughout the table. The cited conductor theorem supplies its exhaustiveness.

\subsection{What the same test gives at exponent seven}
For exact exponent seven, EKN's initial ceiling is $3.9\cdot10^{30}$. Outside their completed census the class number is odd, so genus theory gives one ramified prime. The only possibilities outside the small dyadic cases have $D=-x$ with $x$ prime. In particular there is no ramified prime in the test supports below. If no prime splits up to $Y$, the lower bound for $\log x$ becomes $S+2-14C/9$. We give the intermediate cutoffs to make the implication reproducible.
\begin{center}\begin{tabular}{rrr}\toprule
$Y$ & previous $\log x<$ & forced $\log x>$ if no prime splits\\\midrule
7000 & 70.50 & 76.0\\
5000 & 63.37 & 63.6\\
4700 & 61.01 & 61.5\\
4600 & 60.58 & 60.8\\\bottomrule
\end{tabular}\end{center}
Each contradiction produces a split $p\le Y$, so $x\le 4Y^7$. The final bound is $4\cdot4600^7=1.743270628864\cdot10^{26}$. EKN found 33 fields in the completed range; this calculation leaves the interval between their search and the new upper limit open.

\medskip\noindent\textbf{Exact arithmetic.}
If $v_i$ are the displayed heights, $F(j\delta)=
(\sum_i v_i v_{i+j})/(\sum_i v_i^2)$, taking heights to be zero
outside the twelve intervals. Interpolation between consecutive knots
is linear. The accompanying \texttt{zero\_test.py} compares
$(p^k)^{12}$ with powers of $Y$ to locate the interval exactly. It
encloses logarithms by reducing to $1\le z\le2$ and taking 20 terms
of the series for $2\operatorname{arctanh}((z-1)/(z+1))$, with its
geometric tail. Integer square roots give bounds on $1/\sqrt{p^k}$.
The code then verifies each finite prime sum, each slope bound and
the chain of contradictions using rationals. No zero ordinate or class
group is supplied as numerical input to these tests.
